\documentclass[10pt,a4paper]{amsart}
\usepackage[margin=19mm]{geometry}
\usepackage{amsmath,amssymb,mathtools}
\usepackage[scr=rsfs,cal=euler]{mathalfa}
\usepackage{xfrac}
\usepackage[unicode,hidelinks,bookmarks=false]{hyperref}
\newcommand{\ie}{i.e.\ }
\newcommand{\cf}{cf.\ }
\newcommand{\resp}{respectively\ }
\newcommand{\Subsection}{\subsection}
\providecommand{\nobreakdash}{\nobreak\hbox{-}}
\setlength{\emergencystretch}{2em}
\multlinegap0pt
\usepackage{xcolor}
\usepackage{bm}
\usepackage{mathtools}
\usepackage{algorithm}\usepackage{algpseudocode}
\newtheorem{lemma}{Lemma}
\newtheorem{theorem}{Theorem}
\newcommand{\F}{\mathbf{F}}        % Forman curvature
\newcommand{\Faug}{\F^{\mathrm{aug}}}
\newcommand{\K}{\mathcal{K}}
\title{Forman--Ricci Curvature on Contact-Sequence Temporal Networks via Spatiotemporal Prism Complexes}
\author{Taiki Yamada}
\date{Source: arXiv:2605.15685v1 (CC BY 4.0)}
\begin{document}
\maketitle

\noindent\emph{Source and licence.} Forman--Ricci Curvature on Contact-Sequence Temporal Networks via Spatiotemporal Prism Complexes. Version arXiv:2605.15685v1 (CC BY 4.0), distributed by arXiv under the Creative Commons Attribution 4.0 International licence, http://creativecommons.org/licenses/by/4.0/, as recorded in the arXiv licence metadata for this exact version and shown on the abstract page https://arxiv.org/abs/2605.15685v1. Full licence notice: \texttt{licenses/arxiv-2605.15685-CC-BY-4.0.txt}.\par\smallskip

\noindent\textit{Editorial scope.} Counted unit: the article's theorem prop:monotonicity (source0.tex lines 616--631). The article's complete original proof of it is delivered with it (source0.tex lines 633--692, one \texttt{proof} environment of 60 lines). No mathematics is written, repaired or re-derived: the statement and the proof are the article's own lines, in the article's own notation, and no retained line is substituted or re-flowed.  Retained uncounted as the source-local context the proof needs: lines 265--268; lines 384--394; lines 411--417.  Nothing was omitted from any retained span, so the package carries no omission record.  No retained line contains a citation, so the wrapper carries no bibliography and no bibliography entry.  No retained line contains a figure environment, an image-inclusion command or drawing code, so no image is delivered and the package declares no asset dependency.  Editorial formatting only, in addition: an AMS \texttt{amsart} preamble that restates the article's own declarations and macros used by the retained lines, namely the environments \texttt{lemma}, \texttt{theorem} on separate counters, together with the standard packages those lines need.  The retained environments print in this document as Lemma 1, Theorem 1 in order of appearance, whereas the article prints the same items with the numbering of its own full document.  The complete native source of the article as distributed in the arXiv e-print for this exact version is bundled unchanged as \texttt{sources/200673/source0.tex}.
\begin{equation}\label{eq:prism}
    S_i = \{\hat v_0,\dots,\hat v_i,w_i,w_{i+1},\dots,w_n\},
    \quad i=0,\dots,n,
\end{equation}

\begin{lemma}\label{lem:reduction}
Let $\K$ be an abstract simplicial complex, and let $\alpha,\hat\alpha$ be two distinct $1$-simplices in $\K$.
Then the following equivalence holds:
\begin{equation*}
    \hat\alpha\Vert\alpha
    \iff
    \alpha\cap\hat\alpha\neq\emptyset
    \;\text{and}\;
    \nexists\, \beta\in\K^{(2)}\text{ with }\alpha,\hat\alpha\subset\beta.
\end{equation*}
\end{lemma}

\begin{lemma}\label{lem:agreement}
Let $\K$ be a finite simplicial complex on which all simplex weights equal the same positive constant.
Subsequently, for every $1$-simplex $e\in\K^{(1)}$,
\begin{equation*}
    \F(e) \;=\; \Faug(e).
\end{equation*}
\end{lemma}

\begin{theorem}\label{prop:monotonicity}
Let $\K_{\mathrm{ST}}(\mathcal{C})$ be constructed using a slice gap $K \ge 1$ and uniform weights satisfying $w \equiv 1$.
Let $e^{(k)} = \{(u, t_k), (v, t_k)\}$ denote a spatial edge of $\K_{\mathrm{ST}}$ at time $t_k$.
Assuming that the following conditions hold:
\begin{enumerate}
\item[\textnormal{(A1)}] there exists $\delta \in \{1, \ldots, K\}$ such that $e \in F_{t_{k+s}}$ for every $|s| \le \delta$,
\item[\textnormal{(A2)}] every $\hat e \in \mathcal{P}_{\mathrm{prism}}(e^{(k)})$ is generated by Hatcher's prism operator applied to a persistent simplex that contains both endpoints of $e$.
\end{enumerate}
Then,
\begin{eqnarray}\label{eq:monotonicity_main}
\F_{\K_{\mathrm{ST}}}(e^{(k)}) - \F_{F_{t_k}}(e)
&=& |\mathrm{cof}_2^{\mathrm{prism}}(e^{(k)})| - |\mathcal{P}_{\mathrm{prism}}(e^{(k)})|\nonumber\\
&\ge& 0.
\end{eqnarray}
The equality holds if and only if every $T \in \mathrm{cof}_2^{\mathrm{prism}}(e^{(k)})$ contains a temporal edge as one of its two non-$e^{(k)}$ edges.
\end{theorem}

\begin{proof}
Under uniform weights, Lemma~\ref{lem:agreement} yields the closed-form expression:
$\F(e) = 2 + |\mathrm{cof}_2(e)| - |\mathcal{P}(e)|$ for any $1$-simplex $e$ in the simplicial complex.
Applying this expression to $e^{(k_0)} \in \K_{\mathrm{ST}}$ and $e \in F_{t_{k_0}}$ yields
\begin{align*}
\F_{\K_{\mathrm{ST}}}(e^{(k)}) &= 2 + |\mathrm{cof}_2^{\K_{\mathrm{ST}}}(e^{(k)})| - |\mathcal{P}(e^{(k)})|, \\
\F_{F_{t_k}}(e) &= 2 + |\mathrm{cof}_2^{F_{t_k}}(e)| - |\mathcal{P}_{F_{t_k}}(e)|.
\end{align*}
Therefore,
\begin{eqnarray}\label{eq:result_step1}
\F_{\K_{\mathrm{ST}}}(e^{(k)}) - \F_{F_{t_{k}}}(e) &=& (|\mathrm{cof}_2^{\K_{\mathrm{ST}}}(e^{(k)})| - |\mathrm{cof}_2^{F_{t_k}}(e)|)\nonumber\\
&-& (|\mathcal{P}(e^{(k)})| - |\mathcal{P}_{F_{t_k}}(e)|).
\end{eqnarray}
The cofaces of $e^{(k_0)}$ in $\K_{\mathrm{ST}}$ are categorized into spatial $T \in F_{t_{k_0}}$ and non-spatial $T \in \mathrm{cof}_2^{\mathrm{prism}}(e^{(k_0)})$. Then,
\begin{eqnarray*}
|\mathrm{cof}_2^{\K_{\mathrm{ST}}}(e^{(k)})| &=& \{T \in \K_{\mathrm{ST}}^{(2)} : T \supset e^{(k)},\ T \in F_{t_k}\}\\
&+& \{T \in \K_{\mathrm{ST}}^{(2)} : T \supset e^{(k)},\ T \notin F_{t_k}\}\\
&=& |\mathrm{cof}_2^{F_{t_k}}(e)| + |\mathrm{cof}_2^{\mathrm{prism}}(e^{(k)})|.
\end{eqnarray*}

For the parallel sets, we first observe  that any $T \in \mathrm{cof}_2^{\mathrm{prism}}(e^{(k)})$ has at least one vertex outside the time slice $t_k$ by definition.
By inspecting the top-dimesnional simplices in Equation~\eqref{eq:prism}, we see that any prism-origin 2-simplex contains at most one $1$-simplex whose endpoints both lie in the time slice $t_k$, namely the bottom face of the prism.
Therefore, $T$ cannot simultaneously contain $e^{(k)}$ and another spatial $1$-simplex $\hat e \in F_{t_k}$.
This argument shows that any $\hat e \in F_{t_k} \cap \mathcal{P}_{F_{t_k}}(e)$ remains parallel to $e^{(k)}$ in $\K_{\mathrm{ST}}$, since no new common 2-coface is created. Consequently, the spatial parallel set is in bijective correspondence with $\mathcal{P}_{F_{t_k}}(e)$.
Conversely, assumption (A2) implies that any non-spatial parallel edge $\hat e \in \mathcal{P}(e^{(k)})$ with $\hat e \notin F_{t_k}$ is generated by prism subdivision and therefore belongs to $\mathcal{P}_{\mathrm{prism}}(e^{(k)})$.
Combining these observations, we have
\begin{equation*}
|\mathcal{P}(e^{(k)})| - |\mathcal{P}_{F_{t_k}}(e)| = |\mathcal{P}_{\mathrm{prism}}(e^{(k)})|.
\end{equation*}
Therefore, we obtain Equation ~\eqref{eq:monotonicity_main}.

Further, we show $|\mathcal{P}_{\mathrm{prism}}(e^{(k)})| \le |\mathrm{cof}_2^{\mathrm{prism}}(e^{(k)})|$.
We construct an injection map as $\Phi: \mathcal{P}_{\mathrm{prism}}(e^{(k)}) \to \mathrm{cof}_2^{\mathrm{prism}}(e^{(k)})$.
Let $\hat e \in \mathcal{P}_{\mathrm{prism}}(e^{(k)})$.
By Lemma~\ref{lem:reduction}, $\hat e$ shares exactly one vertex with $e^{(k)}$. Without loss of generality, assuming that this common vertex is $(u, t_k)$, ssuch that $\hat e = \{(u, t_k), (w_*, t_{k+s_*})\}$ for some $w_* \in V$ satisfying $|s_*| \le \delta$.
By assumption (A2), there exists a persistent simplex $\sigma$ such that $\{u, w_*\} \subseteq \sigma$, $\sigma \in F_{t_k} \cap F_{t_{k+s_*}}$. This simplex generates $\hat e$ through prism subdivision. Assumption (A2) further implies that $v \in \sigma$; therefore, $\{u, v, w_*\} \in F_{t_k} \cap F_{t_{k+s_*}}$.
Further, we define
\begin{equation*}
\Phi(\hat e) := \{(u, t_k), (v, t_k), (w_*, t_{k+s_*})\} \in \K_{\mathrm{ST}}^{(2)},
\end{equation*}
whose three edges are $e^{(k)}$, $\hat e$, and $\{(v, t_k), (w_*, t_{k+s_*})\}$, where the last edge arises from prism subdivision applied to $\{u, v, w_*\}$.
Moreover, $\Phi(\hat e) \notin F_{t_k}$, as one of its vertices lies at time $t_{k+s_*} \neq t_k$. Therefore, $\Phi(\hat e) \in \mathrm{cof}_2^{\mathrm{prism}}(e^{(k)})$.

To demonstrate that $\Phi$ is injective, we assume that $\Phi(\hat e_1) = \Phi(\hat e_2) = T$.
The simplex $T$ has exactly three edges, two of which are $e^{(k)}$ and the corresponding preimage edge by construction, leaving only one remaining edge, denoted by $f$.
If $\hat e_1$ and $\hat e_2$ are distinct preimages, then one of them can coincide with $f$. However, $f$ shares a vertex with $e^{(k)}$, and both edges lie in the common $2$-coface $T$.
From Lemma~\ref{lem:reduction}, we have $f \notin \mathcal{P}(e^{(k)})$, which contradicts the assumption that $f = \hat e_i \in \mathcal{P}(e^{(k)})$.
Hence, $\hat e_1 = \hat e_2$.

Therefore, the injection map $\Phi$ yields
\begin{equation*}
|\mathcal{P}_{\mathrm{prism}}(e^{(k)})| \le |\mathrm{cof}_2^{\mathrm{prism}}(e^{(k)})|,
\end{equation*}
thereby completing Equation ~\eqref{eq:monotonicity_main}.

Furthermore, $\Phi$ is a bijection if and only if every $T \in \mathrm{cof}_2^{\mathrm{prism}}(e^{(k)})$ lies in the image.
By the preceding argument, among the two non-$e^{(k)}$ edges of $T$, at least one fails to be parallel to $e^{(k)}$ because it shares the common $2$-coface $T$.
Hence, the only way for the remaining edge to be parallel to $e^{(k)}$ and to lie in the image of $\Phi$ is for it to be a temporal edge $\{(v, t_k), (v, t_{k+s_*})\}$ (or the analogous edge associated with $u$) generated solely by the persistence of $\{v\}$.
This condition is precisely the equality condition stated in the theorem.
\end{proof}

\end{document}
